%%% Внесите свои данные - Input your data
%%
%%
\newcommand{\Author}{Н.Ю.\,Каралюс} % И.О. Фамилия автора
\newcommand{\AuthorFull}{Каралюс Никита Юргисович} % Фамилия Имя Отчество автора
\newcommand{\AuthorFullDat}{Каралюсу Никите Юргисовичу} % Фамилия Имя Отчество автора в дательном падеже (Кому? Студенту...)
\newcommand{\AuthorFullVin}{Каралюса Никиты Юргисовича} % в винительном падеже (Кого? что?  Програмиста ...)
\newcommand{\AuthorPhone}{+7-993-533-22-98} % номер телефорна автора для оперативной связи
\newcommand{\Supervisor}{А.В.\,Сергеев} % И. О. Фамилия научного руководителя
\newcommand{\SupervisorFull}{Сергеев Анатолий Васильевич} % Фамилия Имя Отчество научного руководителя
\newcommand{\SupervisorVin}{А.В.\,Сергеева} % И. О. Фамилия научного руководителя  в винительном падеже (Кого? что? Руководителя ...)
\newcommand{\SupervisorJob}{доцент ВШ ПИ} %
\newcommand{\SupervisorJobVin}{доцента ВШ ПИ} % в винительном падеже (Кого? что?  Програмиста ...)
\newcommand{\SupervisorDegree}{к.т.н} %
\newcommand{\SupervisorTitle}{} %
%%
%%
%Руководитель, УТВЕРЖДАЮЩИЙ задание
\newcommand{\Head}{А.В.\,Щукин} % И. О. Фамилия руководителя подразделения (руководителя ОП)
\newcommand{\HeadDegree}{Руководитель ОП}% Только должность:
%Руководитель ОП % <- ПРАВИЛЬНЫЙ ВАРИАНТ ДЛЯ 09.03.03 и 09.04.03
%Заведующий % кафедрой
%Директор % Высшей школы
%Зам. директора
\newcommand{\HeadDep}{} % заменить на краткую аббревиатуру подразделения или ОСТАВИТЬ ПУСТЫМ, если утверждает руководитель ОП

%%% Руководитель, принимающий заявление
\newcommand{\HeadAp}{А.В.\,Щукин} % И. О. Фамилия руководителя подразделения (руководителя ОП)
\newcommand{\HeadApDegree}{Руководитель ОП}% Только должность:
%Руководитель ОП
%Заведующий кафедрой
%Директор Высшей школы
\newcommand{\HeadApDep}{} % заменить на краткую аббревиатуру подразделения или оставить пустым, если утверждает руководитель ОП
%%% Консультант по нормоконтролю
\newcommand{\ConsultantNorm}{Е.Е.\,Андрианова} % И. О. Фамилия консультанта по нормоконтролю. ТОЛЬКО из числа ППС!
\newcommand{\ConsultantNormDegree}{ст.преподаватель ВШ ПИ} %
%%% Первый консультант
\newcommand{\ConsultantExtraFull}{Маркелов Максим Александрович} % Фамилия Имя Отчетство дополнительного консультанта
\newcommand{\ConsultantExtra}{М.А.\,Маркелов} % И. О. Фамилия дополнительного консультанта
\newcommand{\ConsultantExtraDegree}{тех.руководитель команды разработки} %
\newcommand{\ConsultantExtraVin}{М.А.\,Маркелова} % И. О. Фамилия дополнительного консультанта в винительном падеже (Кого? что? Руководителя ...)
\newcommand{\ConsultantExtraDegreeVin}{тех.руководителя команды разработки} %  в винительном падеже (Кого? что? Руководителя ...)
%%% Второй консультант
\newcommand{\ConsultantExtraTwoFull}{Фамилия Имя Отчетство} % Фамилия Имя Отчетство дополнительного консультанта
\newcommand{\ConsultantExtraTwo}{И.О.\,Фамилия} % И. О. Фамилия дополнительного консультанта
\newcommand{\ConsultantExtraTwoDegree}{должность, степень} %
\newcommand{\ConsultantExtraTwoVin}{И.О.\,Фамилию} % И. О. Фамилия дополнительного консультанта в винительном падеже (Кого? что? Руководителя ...)
\newcommand{\ConsultantExtraTwoDegreeVin}{должность, степень} %  в винительном падеже (Кого? что? Руководителя ...)
\newcommand{\Reviewer}{И.О.\,Фамилия} % И. О. Фамилия резензента. Обязателен только для магистров.
\newcommand{\ReviewerDegree}{должность, степень} %
%%
%%
\renewcommand{\thesisTitle}{Система автоматической бесшовной ротации криптографических ключей авторизации в высоконагруженной микросервисной платформе (на примере ООО~<<Майндбокс>>)}
\newcommand{\thesisDegree}{работа бакалавра}% магистерская диссертация %c 2020, И ТОЛЬКО для специалитата дипломный проект и дипломная работа
\newcommand{\thesisTitleEn}{Automated zero-downtime cryptographic key rotation system for authorization in a large-scale microservice platform (a case study of Mindbox LLC)} %2020
\newcommand{\thesisDeadline}{18.05.2026} %Последний день преддипломной практики согласно учебному плану.
\newcommand{\thesisStartDate}{16.02.2026}
\newcommand{\thesisYear}{2026} % заменить на год защиты
\newcommand{\approveYear}{202\underline{\hspace*{0.01\textheight}}} % <- НЕ МЕНЯТЬ, ТОЛЬКО С 2030го :)
%%
%%
\newcommand{\group}{5130903/20302} % заменить вместо N номер группы
\newcommand{\thesisSpecialtyCode}{09.03.03}% код направления подготовки
\newcommand{\thesisSpecialtyTitle}{Прикладная информатика} % наименование направления/специальности
\newcommand{\thesisOPPostfix}{03} % последние цифры кода образовательной программы (после <<_>>)
\newcommand{\thesisOPTitle}{Интеллектуальные инфокоммуникационные технологии}% наименование образовательной программы

\newcommand{\institute}{Институт компьютерных наук и~кибербезопасности}




%%% Задание ключевых слов и аннотации
%%
%%
%% Ключевых слов от 3 до 5 слов или словосочетаний в именительном падеже именительном падеже множественного числа (или в единственном числе, если нет другой формы) по правилам русского языка!!!
%%
%%
\newcommand{\keywordsRu}{безопасность, JWT, межсервисная авторизация, ротация ключей, микросервисы} % ВВЕДИТЕ ключевые слова по-русски
%%
%%
\newcommand{\keywordsEn}{security, JWT, service-to-service authorization, key rotation, microservices} % ВВЕДИТЕ ключевые слова по-английски
%%
%%
%% Реферат ОТ 1000 ДО 1500 знаков на русский или английский текст
%%
%Реферат должен содержать:
%- предмет, тему, цель ВКР;
%- метод или методологию проведения ВКР:
%- результаты ВКР:
%- область применения результатов ВКР;
%- выводы.

\newcommand{\abstractRu}{Объект исследования~--- система межсервисной авторизации высоконагруженной микросервисной платформы Mindbox (около 160~микросервисов, свыше 2~млн запросов в минуту). Цель работы~--- спроектировать, реализовать и внедрить систему автоматической бесшовной ротации криптографических ключей подписи JWT. Методологической основой работы выступает системный анализ: декомпозиция системы межсервисной авторизации на взаимосвязанные элементы, анализ их связей и ограничений, синтез целевой архитектуры и её экспериментальная апробация. Разработана система из пяти компонентов. Бесшовность обеспечивает трёхшаговый алгоритм ротации с grace-периодами и проверкой доступности ключей; межконтурное взаимодействие построено на Token Exchange. Использованы C\#, Go, Docker, Kubernetes, PostgreSQL и Yandex Cloud. Система внедрена в промышленную эксплуатацию: ключи ротируются автоматически по расписанию, время ответа по 95-му процентилю~--- 8~мс при SLO 30~мс, подтверждённая пропускная способность~--- 15~000 запросов в секунду, за два месяца эксплуатации инцидентов авторизации не зафиксировано. Результаты применимы в высоконагруженных микросервисных платформах, требующих плановой смены ключей без простоя. Подтверждена гипотеза: бесшовная автоматическая ротация достижима собственным решением на основе открытых стандартов с учётом особенностей существующей архитектуры.} % ВВЕДИТЕ текст аннотации по-русски
%%
%%
\newcommand{\abstractEn}{This work addresses service-to-service authorization in Mindbox, a large-scale microservice platform that runs roughly 160~services and handles over 2~million requests per minute. The aim is to design, implement, and deploy a system that automatically rotates cryptographic JWT signing keys with zero downtime. Methodologically, the work rests on systems analysis: it decomposes the service-to-service authorization system into interrelated elements, analyzes their relationships and constraints, synthesizes a target architecture, and validates it experimentally. The resulting system comprises five components. A three-step rotation algorithm with grace periods and key-availability checks delivers zero downtime, while cross-environment requests rely on Token Exchange. The system is written in C\# and Go, packaged with Docker and Kubernetes, backed by PostgreSQL, and deployed on Yandex Cloud. It now runs in production: keys rotate automatically on a schedule, the 95th-percentile response time is 8~ms against a 30~ms SLO, measured throughput reaches 15,000~requests per second, and no authorization incidents have occurred in two months of operation. The results apply to any large-scale microservice platform that needs scheduled, zero-downtime key rotation, and they confirm the hypothesis: such rotation is achievable with an in-house solution built on open standards and tailored to the platform's existing architecture.} % ВВЕДИТЕ текст аннотации по-английски


%%% РАЗДЕЛ ДЛЯ ОФОРМЛЕНИЯ ПРАКТИКИ
%Место прохождения практики
\newcommand{\PracticeType}{Отчет о прохождении %
	%стационарной производственной (технологической (проектно-технологической)) %
	такой-то % тип и вид ЗАМЕНИТЬ
	практики}

%%% Руководитель практической подготовки от профильной организации
\newcommand{\PracticeOrgSupervisorFull}{Маркелов Максим Александрович}
\newcommand{\PracticeOrgSupervisor}{М.А.\,Маркелов}
\newcommand{\PracticeOrgSupervisorDegree}{тех.руководитель команды разработки}

\newcommand{\Workplace}{ФГАОУ ВО «СПбПУ», ИКНК, ВШ ПИ}

% Даты начала/окончания
\newcommand{\PracticeStartDate}{%
01.09.2025%
%	22.06.2020
}%
\newcommand{\PracticeEndDate}{%
16.12.2025%
%	18.07.2020%
}%
%%

\newcommand{\School}{
Высшая школа программной инженерии
}
\newcommand{\practiceTitle}{Система автоматической бесшовной ротации криптографических ключей авторизации в высоконагруженной микросервисной платформе (на примере ООО~<<Майндбокс>>)}


%% ВНИМАНИЕ! Необходимо либо заменить текст аннотации (ключевых слов) на русском и английском, либо удалить там весь текст, иначе в свойства pdf-отчета по практике пойдет шаблонный текст.

%%% Не меняем дальнейшую часть - Do not modify the rest part
%%
%%
%%
%%
\ifnumequal{\value{docType}}{1}{% Если ВКР, то...
	\newcommand{\DocType}{Выпускная квалификационная работа}
	\newcommand{\pdfDocType}{\DocType~(\thesisDegree)} %задаём метаданные pdf файла
	\newcommand{\pdfTitle}{\thesisTitle}
}{% Иначе
	\newcommand{\DocType}{\PracticeType}
	\newcommand{\pdfDocType}{\DocType} %задаём метаданные pdf файла
	\newcommand{\pdfTitle}{\practiceTitle}
}%
\newcommand{\HeadTitle}{\HeadDegree~\HeadDep}
\newcommand{\HeadApTitle}{\HeadApDegree~\HeadApDep}
\newcommand{\thesisOPCode}{\thesisSpecialtyCode\_\thesisOPPostfix}% код образовательной программы
\newcommand{\thesisSpecialtyCodeAndTitle}{\thesisSpecialtyCode~\thesisSpecialtyTitle}% Код и наименование направления/специальности
\newcommand{\thesisOPCodeAndTitle}{\thesisOPCode~\thesisOPTitle} % код и наименование образовательной программы
%%
%%
\hypersetup{%часть болка hypesetup в style
		pdftitle={\pdfTitle},    % Заголовок pdf-файла
		pdfauthor={\AuthorFull},    % Автор
		pdfsubject={\pdfDocType. Шифр и наименование направления подготовки: \thesisSpecialtyCodeAndTitle. \abstractRu},      % Тема
		pdfcreator={LaTeX, SPbPU-student-thesis-template},     % Приложение-создатель
%		pdfproducer={},  % Производитель, Производитель PDF % будет выставлена автоматически
		pdfkeywords={\keywordsRu}
}
%%
%%
%% вспомогательные команды
\newcommand{\firef}[1]{рис.\ref{#1}} %figure reference
\newcommand{\taref}[1]{табл.\ref{#1}}	%table reference
%%
%%
%% Архивный вариант задания ключевых слов, аннотации и благодарностей
% Too hard to export data from the environment to pdf-info
% https://tex.stackexchange.com/questions/184503/collecting-contents-of-environment-and-store-them-for-later-retrieval
%заменить NewEnviron на newenvironment для распознавания команды в TexStudio
%\NewEnviron{keywordsRu}{\noindent\MakeUppercase{\BODY}}
%\NewEnviron{keywordsEn}{\noindent\MakeUppercase{\BODY}}
%\newenvironment{abstractRu}{}{}
%\newenvironment{abstractEn}{}{}
%\newenvironment{acknowledgementsRu}{\par{\normalfont \acknowledgements.}}{}
%\newenvironment{acknowledgementsEn}{\par{\normalfont \acknowledgementsENG.}}{}


%%% Переопределение именований %%% Не меняем - Do not modify
%\newcommand{\Ministry}{Минобрнауки России}
\newcommand{\Ministry}{Министерство науки и высшего образования Российской~Федерации} %с 2020
\newcommand{\SPbPU}{Санкт-Петербургский политехнический университет Петра~Великого}
\newcommand{\SPbPUOfficialPrefix}{Федеральное государственное автономное образовательное учреждение высшего образования}
\newcommand{\SPbPUOfficialShort}{ФГАОУ~ВО~<<СПбПУ>>}
%% Пробел между И. О. не допускается.
\renewcommand{\alsoname}{см. также}
\renewcommand{\seename}{см.}
\renewcommand{\headtoname}{вх.}
\renewcommand{\ccname}{исх.}
\renewcommand{\enclname}{вкл.}
\renewcommand{\pagename}{Pages}
\renewcommand{\partname}{Часть}
\renewcommand{\abstractname}{\textbf{Аннотация}}
\newcommand{\abstractnameENG}{\textbf{Annotation}}
\newcommand{\keywords}{\textbf{Ключевые слова}}
\newcommand{\keywordsENG}{\textbf{Keywords}}
\newcommand{\acknowledgements}{\textbf{Благодарности}}
\newcommand{\acknowledgementsENG}{\textbf{Acknowledgements}}
\renewcommand{\contentsname}{Content} %
%\renewcommand{\contentsname}{Содержание} % (ГОСТ Р 7.0.11-2011, 4)
%\renewcommand{\contentsname}{Оглавление} % (ГОСТ Р 7.0.11-2011, 4)
\renewcommand{\figurename}{Рис.} % Стиль СПбПУ
%\renewcommand{\figurename}{Рисунок} % (ГОСТ Р 7.0.11-2011, 5.3.9)
\renewcommand{\tablename}{Таблица} % (ГОСТ Р 7.0.11-2011, 5.3.10)
%\renewcommand{\indexname}{Предметный указатель}
\renewcommand{\listfigurename}{Список рисунков}
\renewcommand{\listtablename}{Список таблиц}
\renewcommand{\refname}{\fullbibtitle}
\renewcommand{\bibname}{\fullbibtitle}

\newcommand{\chapterEnTitle}{Сhapter title} % <- input the English title here (only once!)
\newcommand{\chapterRuTitle}{Название главы}          % <- введите
\newcommand{\sectionEnTitle}{Section title} %<- input subparagraph title in english
\newcommand{\sectionRuTitle}{Название подраздела} % <- введите название подраздела по-русски
\newcommand{\subsectionEnTitle}{Subsection title} % - input subsection title in english
\newcommand{\subsectionRuTitle}{Название параграфа} % <- введите название параграфа по-русски
\newcommand{\subsubsectionEnTitle}{Subsubsection title} % <- input subparagraph title in english
\newcommand{\subsubsectionRuTitle}{Название подпараграфа} % <- введите название подпараграфа по-русски
